\chapter{Future Work}
\label{ch:future-work}

This chapter lists directions in which the work on the Cyclops rover can be extended, first for
monocular depth estimation and then for localization.

\begin{itemize}
  \item \textbf{Depth Anything V3.} Depth Anything V3~\red{[]} was released after the depth
  pipeline of Chapter~\ref{ch:mde} was built around DAv2~\cite{DepthAnythingV2}. Replacing the
  backbone with the newer model and repeating the accuracy and latency tests of
  Chapter~\ref{ch:results} would show whether it gives a better range scan on the rover.

  \item \textbf{Metric models versus calibrated relative depth.} The metric DAv2 models were set
  aside in Section~\ref{sec:depth-calibration} on the basis of preliminary trials. A direct
  comparison is still needed: the metric models and our relative model, after the scale and
  shift calibration of Section~\ref{sec:metric-scale}, should be evaluated on the same scenes to
  see which gives more accurate depth on the rover.

  \item \textbf{Ground points from mounting geometry.} The ground points used to calibrate depth
  are currently found by placing a checkerboard on the floor and solving the
  Perspective-n-Point problem (Section~\ref{sec:metric-scale}). Since the camera is rigidly
  mounted, the 3D coordinate of any floor pixel can instead be computed directly from the
  camera's height and tilt, by intersecting the pixel's ray with the ground plane. This removes
  the checkerboard altogether and allows ground points closer to the rover, which would reduce
  the minimum range discussed in Section~\ref{sec:results-limitations}.

  \item \textbf{Depth accuracy away from the floor.} The depth accuracy test of
  Section~\ref{sec:results-depth} measures error only at checkerboard corners, which all lie on
  the floor. The obstacles that form the range scan, however, are walls and objects above the
  floor. Recording the same scenes with an RGB-D camera~\cite{rgbd} or a LiDAR as ground truth
  would measure the depth error at these points as well.

  \item \textbf{Learning the measurement model parameters.} The beam model parameters
  $\sigma_{\mathrm{hit}}$, $\lambda_{\mathrm{short}}$, and the mixture weights $\beta$
  (Table~\ref{tab:meas-params}) can be learned from data instead of being set by hand. Thrun et
  al.~\cite{thrun2005probabilistic} describe a maximum likelihood estimator that finds the
  parameters that best explain a set of recorded scans taken from known poses. Applying it to
  MDE range scans would fit the model to the actual error of the depth network.

  \item \textbf{Simultaneous localization and mapping.} The rover currently localizes against a
  map that is made in advance. A natural next step is simultaneous localization and mapping
  (SLAM)~\cite{durrantwhyte2006slam}, in which the rover builds the occupancy grid from its own
  MDE range scans while localizing in it, so that no prior map is needed.
\end{itemize}
%%% Local Variables: ***
%%% mode: latex ***
%%% TeX-master: "thesis.tex" ***
%%% End: ***
